\documentclass{article}
\usepackage{amsmath,amsthm,amssymb}
\usepackage{array}
\usepackage{url}

\title{Finite Computations in the Dyck Symmetric Functions Appendices}
\author{}
\date{}

\begin{document}
\emergencystretch=3em
\maketitle

\section{What the Appendix Code Supports}

The 2026 preprint proves a string decomposition for Dyck sequences in a range
of deficits. The construction uses two local operations, called
\(\mathrm{up}\) and \(\mathrm{down}\), which move a Dyck sequence one step in
area while preserving the relevant deficit. These operations are defined by
looking for a small local pattern, removing a few entries, changing a short
East or West window, and then reinserting the changed entries.

Most of the proof is symbolic: it shows that the operations preserve the Dyck
condition, preserve the deficit, and stay inside the intended range. The
appendix computations handle finite parts of this argument where the proof has
reduced the problem to checking all possible short windows or all possible
small residual inputs.

\section{Appendix A}

Appendix A gives reusable code for Dyck sequences. It defines the basic
statistics, the skeleton tests, the insertion and extraction routines, and the
local \(\mathrm{up}\) and \(\mathrm{down}\) maps. It also gives the routine
that starts from the special skeletons and repeatedly applies
\(\mathrm{up}\) to form the lower half of each string.

The Appendix A code is included here because it is the computational form of
the construction itself. It can be used to reproduce examples of the strings
and to check small cases of the decomposition.

In this repository, Appendix A is represented by the following files in
\path{code/appendix_a/}:
\begin{itemize}
\item \path{01_core_dyck_sequence_routines.py}: the
basic Dyck-sequence routines and the definitions of the local
\(\mathrm{up}\) and \(\mathrm{down}\) algorithms.
\item \path{02_make_strings.py}: the Appendix A routine
that builds the lower-half strings from the special skeletons.
\end{itemize}

\section{Appendix B}

Appendix B proves that the local operations used in Appendix A are
well-defined in the range needed by the theorem. In the arXiv version, the
four local lemmas are Lemmas 5.22--5.25:
\begin{center}
\begin{tabular}{>{\raggedright\arraybackslash}p{0.18\linewidth}
                >{\raggedright\arraybackslash}p{0.34\linewidth}
                >{\raggedright\arraybackslash}p{0.34\linewidth}}
Lemma & Title & Role \\
\hline
5.22 & Skeleton cases succeed & Checks that the skeleton branches return Dyck
sequences in the required range. \\
5.23 & Extraction chains never fail & Checks that every extraction requested
by the staged algorithms exists. \\
5.24 & The seven-window branches do not fail & Checks the final seven-entry
East and West local moves. \\
5.25 & Bounded extraction positions and injection nonfailure & Checks the
position bounds needed for later insertion steps. \\
\end{tabular}
\end{center}

These lemmas check four basic points:
\begin{itemize}
\item when the construction recognizes a skeleton case, the proposed output is
again a Dyck sequence of the required form;
\item when the construction needs to remove entries, the entries it asks for
actually exist;
\item when the construction reaches the seven-entry East or West move, every
possible short window satisfies the inequalities needed by the proof;
\item the positions used for later insertion steps stay within the allowed
bounds.
\end{itemize}

The symbolic proof handles the infinite families. The finite computations in
Appendix B close the remaining cases.

\subsection{Residual Small Range}

One checker enumerates the remaining small values \(4\le n\le 7\). For each
Dyck sequence in that finite range, it follows the decision steps used by the
\(\mathrm{up}\) and \(\mathrm{down}\) algorithms and confirms that each input
falls into one of the cases already covered by the proof. It also confirms that
the seven-entry East and West moves are not reached in this small residual
range.

This corresponds to \path{01_residual_finite_check.py} in
\path{code/appendix_b/}. It supports the small-range parts of Lemmas 5.22,
5.23, 5.24, and 5.25 at once:
\begin{itemize}
\item for Lemma 5.22, it checks the skeleton branches in the residual range;
\item for Lemma 5.23, it checks that the extractions used before each branch
stops exist;
\item for Lemma 5.24, it confirms that the seven-window branch is not reached
for \(4\le n\le 7\);
\item for Lemma 5.25, it checks the relevant position bounds before the
residual branch stops.
\end{itemize}
The successful-output transcript is
\path{02_residual_successful_output.txt}.

\subsection{Seven-Entry East and West Windows}

The largest local move changes a window of seven entries. The proof reduces
this part to a finite list of possible East and West windows together with the
short extensions that can appear next to them. The checker enumerates those
windows and verifies the threshold inequalities used in the argument. This is
needed because the local move has several boundary cases that are easier and
less error-prone to exhaust by computation than to list manually in the text.

This corresponds to \path{03_east7_west7_seven_window_checker.py} in
\path{code/appendix_b/}. It is the finite checker for Lemma 5.24 in the
\(n\ge 8\) part of the proof. The successful-output transcript is
\path{04_east7_west7_successful_output.txt}.

\subsection{Position Bounds}

Another part of the proof needs to know that certain extraction and insertion
positions remain legal. Appendix B treats the general case symbolically and
then leaves two finite domains to check. The first finite checker covers words
with at most seven nonzero entries in the required range. The second covers two
specific prefix forms that are excluded from the symbolic argument. In both
cases, the code enumerates the finite domain and verifies that every word leads
to one of the contradictions or bounds used in the written proof.

These files correspond to Lemma 5.25:
\begin{itemize}
\item \path{05_lemma_525_limited_nonzero_checker.py}
checks the enlarged finite domain with \(4\le n\le 13\) and at most seven
nonzero entries. Its successful-output transcript is
\path{06_lemma_525_limited_nonzero_successful_output.txt}.
\item \path{07_lemma_525_prefix_checker.py} checks the
two excluded prefix forms in the finite range \(9\le n\le 16\). Its
successful-output transcript is
\path{08_lemma_525_prefix_successful_output.txt}.
\end{itemize}

\section{Why These Checks Are Included}

The computations are included because they are part of the proof, not just
evidence for the theorem. Each computation corresponds to a finite exhaustive
step that appears after the main argument has reduced an infinite statement to
a bounded list of cases. The appendix code records exactly what was checked,
and the successful-output listings record the expected result of those checks.
\end{document}
